\begin{abstract}
Low-light enhancement is usually judged by RGB fidelity, yet a model can improve PSNR while learning an exposure field with the wrong spatial structure. This paper studies that failure mode through \emph{gauge-identifiable exposure-field learning}. We show that local exposure fields are ambiguous up to an additive image-level gauge, and decompose them into a gauge-free spatial shape $S$ and an absolute gauge $\mu$. Based on this decomposition, we introduce centered exposure-shape calibration and UEFB-G, a gauge-controlled evaluation protocol that separates RGB fidelity, global gauge error, and local exposure-shape error. In a full frozen-evidence audit over 500 UEFB-v2 images, 100 LOL-v2-real images, and 100 LOL-v2-synthetic images, centered E-shape calibration improves UEFB gauge-free shape correlation over joint training by +0.1581 (95\% CI [0.1369, 0.1791], FDR $q=2.18\times10^{-43}$), while also improving real and synthetic PSNR by +0.6963 dB and +0.6321 dB. We further show a PSNR-only misranking: an A-gate variant obtains higher UEFB PSNR but negative exposure-shape correlation, whereas the centered E-shape model yields physically coherent fields. The resulting contribution is not a SOTA restoration claim; Retinexformer remains stronger in RGB fidelity. Instead, GIR-Field reframes uneven low-light enhancement as a problem of gauge-identifiable field learning, benchmark design, and claim-calibrated failure analysis.
\end{abstract}
